\documentclass[tikz,border=4pt]{standalone}
\usepackage{amsmath}
\usepackage{tikz-cd}
\begin{document}
% Chapter dependency graph: a single chain, Chapter 0 through 14, with the
% two checkpoint projects as waypoints (diamond render) after Chapters 6
% and 12, plus two dashed skip edges for the named paths that actually
% bypass material (as opposed to the two paths that just change how a
% chapter is read). Source: learning-paths.md.
\begin{tikzcd}[column sep=huge, row sep=small]
\text{0. Setup} \arrow[d] \\
\text{1. Basics} \arrow[d] \\
\text{2. Terminology \& CoC} \arrow[d] \\
\text{3. Functions \& structures} \arrow[d] \\
\text{4. Propositions \& proofs} \arrow[d] \\
\text{5. Tactics} \arrow[d] \\
\text{6. Rigor check} \arrow[d] \\
\langle\text{Checkpoint: Monoid}\rangle \arrow[d] \\
\text{7. Groups} \arrow[d] \\
\text{8. Group theorems} \arrow[d] \\
\text{9. Rings} \arrow[d] \\
\text{10. Ring theorems} \arrow[d] \\
\text{11. Modules} \arrow[d] \\
\text{12. Path algebras} \arrow[d] \\
\langle\text{Checkpoint: Path.length}\rangle \arrow[d] \\
\text{13. Working efficiently} \arrow[d] \\
\text{14. Next steps}
\arrow["\text{already know Lean}"{sloped}, dashed,
  bend right=80, looseness=3, from=2-1, to=4-1]
\arrow["\text{already know Lean}"{sloped}, dashed,
  bend right=70, looseness=2.5, from=6-1, to=9-1]
\arrow["\text{fastest path}"{sloped}, dashed,
  bend left=55, looseness=1.8, from=1-1, to=9-1]
\end{tikzcd}
\end{document}
